La missió BepiColombo explorarà Mercuri i posarà en òrbita el satèl·lit \textit{Mercury Planetary Orbiter} (\textit{MPO}) al voltant del planeta amb un radi orbital mitjà de 3.360~km.

\begin{apartat}{1,25}
Considereu un satèl·lit que fa una òrbita circular al voltant de Mercuri. Deduïu l'expressió de la velocitat orbital del satèl·lit en funció del radi orbital i la massa de Mercuri (indiqueu clarament en quins principis o lleis físiques us baseu per fer la vostra deducció). Amb aquesta expressió, calculeu la velocitat orbital del satèl·lit \textit{MPO} mentre orbita Mercuri. Calculeu quantes voltes haurà fet el satèl·lit \textit{MPO} al planeta Mercuri al cap d'un any terrestre.
\end{apartat}

\begin{apartat}{1,25}
A partir de l'expressió general de l'energia mecànica, obtingueu-ne l'equació per al cas particular d'un satèl·lit en òrbita circular (cal que l'equació final només estigui expressada en funció de $G$, el radi orbital i les masses del satèl·lit i del planeta). Una vegada el satèl·lit \textit{MPO} està orbitant Mercuri, encara té combustible per poder fer maniobres. Considereu que el combustible disponible pot proporcionar una energia de $4,5\cdot 10^{9}\ \mathrm{J}$. Determineu el valor màxim que podria tenir la massa del \textit{MPO} perquè amb l'energia disponible pogués escapar del camp gravitatori de Mercuri.
\end{apartat}

\dades{%
  $G = 6,67\cdot 10^{-11}\ \mathrm{N\,m^2\,kg^{-2}}$\\
  Massa de Mercuri: $M_\mathrm{M} = 3,285\cdot 10^{23}\ \mathrm{kg}$\\
  Any terrestre $= 365,25$ dies}
